\section{División de trabajo entre clasificador y generador}
\label{sec:disc_division}
% Extensión estimada: 2 páginas, cinco párrafos (ficha 6.3).
% Fuentes: sec:res_gen_fuentes, sec:res_gen_ablaciones, sec:res_gen_pipeline (tab:res_gen_pipeline),
% sec:res_gen_comparacion, sec:res_gen_control_mini, sec:res_gen_rag, tab:res_gen_aportes, sec:res_gen_estado_arte.
% Integra el párrafo 4 del antiguo 4.6 y los párrafos 3 y 4 del antiguo 5.6.

Como el diagnóstico no puede salir del generador, el \emph{pipeline} lo toma de Swin-Tiny, que se lo entrega al
generador a través de la cascada (Sección~\ref{sec:met_cascada}). Los resultados del Capítulo~\ref{cp:res_gen}
muestran hasta qué punto esa decisión reparte las tareas entre los dos componentes. Con Mini-VLM y SkinVL,
la consistencia diagnóstica de la cascada coincide con la exactitud de la fuente, tanto con Swin-Tiny, 0.847, como con
DermLIP sin ajuste, 0.716, y también en cada clase (Secciones~\ref{sec:res_gen_fuentes} y~\ref{sec:res_gen_clases}).
Sin la imagen, el diagnóstico afirmado casi no cambia, y con un diagnóstico equivocado el reporte casi siempre lo afirma
en lugar de corregirlo (Sección~\ref{sec:res_gen_ablaciones}). El clasificador aporta, entonces, el diagnóstico, y el
generador, la redacción y la descripción de cada lesión, que sí varía con la imagen.

Este reparto tiene una consecuencia para el diseño de los sistemas que generan reportes médicos: en una cascada, la
fidelidad diagnóstica del reporte se mejora mejorando el clasificador, y no el generador. La comparación entre fuentes
lo muestra en los dos sentidos. Pasar de DermLIP sin ajuste a Swin-Tiny agrega 0.132 de consistencia con el mismo
generador, mientras que cambiar de generador con la misma fuente no la lleva más allá de la exactitud del
clasificador. La consecuencia es también una advertencia para la evaluación. Un sistema que integra un clasificador y
un generador y se evalúa como un todo, como los de la literatura revisada (Sección~\ref{sec:res_gen_estado_arte}),
puede atribuir al generador una fidelidad que proviene del clasificador. En este trabajo, en cambio, el aporte de cada
componente se mide por separado (Tabla~\ref{tab:res_gen_aportes} del Capítulo~\ref{cp:res_gen}).

El ruteo por concordancia agrega al reparto una tercera tarea, la de decidir si el diagnóstico es lo bastante
confiable para afirmarlo. El diseño de la Sección~\ref{sec:met_pipeline} separa así qué diagnóstico se afirma, que
decide el clasificador, de si corresponde afirmarlo, que decide la concordancia con el \gls{knn}. Sobre las 1\,503
imágenes, la consistencia de Mini-VLM pasa de 0.849 con la cascada simple a 0.870 con el \emph{pipeline}, aun con las
derivaciones contadas como casos sin diagnóstico. La línea base que escribe los mismos diagnósticos sin generador
obtiene una consistencia equivalente, de modo que la ganancia puede atribuirse al ruteo, y el generador no agrega
fidelidad, pero tampoco la pierde. Con LLaVA-Med, en cambio, el \emph{pipeline} queda por debajo de esa línea base,
porque en el nivel diferencial el generador deja los candidatos dentro de un razonamiento inconcluso. Es decir, el
generador no suma fidelidad, pero puede restarla, y por eso su elección importa.

Los resultados respaldan, además, las decisiones de diseño del Capítulo~\ref{cp:metodos}. El ancla en Swin-Tiny se
eligió por su F1 macro en validación (Sección~\ref{sec:met_pipeline}), y la cascada la confirma, porque con DermLIP
sin ajuste como fuente los reportes pierden más de una décima de consistencia. Mini-VLM es el generador del
\emph{pipeline} porque transcribe el diagnóstico como el generador especializado, describe tanto como el biomédico más
grande y es el único cuyos pesos entran en una GPU de consumo de 12~GB (Sección~\ref{sec:res_gen_comparacion}). La
variante de control, que ajusta también su modelo de lenguaje, no mejora la fidelidad y deja reportes vacíos
(Sección~\ref{sec:res_gen_control_mini}), de modo que la configuración principal se mantiene. Por último, la
recuperación queda fuera del ruteo como alternativa al clasificador, porque sola no alcanza al voto del \gls{knn} y
junto al clasificador aporta poco (Sección~\ref{sec:res_gen_rag}).

El mismo reparto que hace transparente la fidelidad del reporte tiene un costo. Como el generador transcribe el
diagnóstico y adapta la descripción a él, un error del clasificador llega al reporte acompañado de rasgos que lo hacen
verosímil, como el melanoma descrito como una lesión de pigmentación homogénea de la
Sección~\ref{sec:res_gen_cualitativos}. La calidad de la redacción, por lo tanto, no indica si el diagnóstico es
correcto, y lo que las métricas permiten afirmar sobre la utilidad clínica del reporte es el tema de la sección
siguiente.
